\section*{Chapter \ref{ch:g3:continents-and-oceans} --- The Globe: Continents and Oceans}

\begin{solution}{exo:g3:continents-and-oceans:1}
Asia, Africa, North America, South America, Antarctica, Europe and
Australia. The largest is Asia; the smallest is Australia.
\end{solution}

\begin{solution}{exo:g3:continents-and-oceans:2}
The Pacific, the Atlantic, the Indian, the Southern and the Arctic \omterm{def:g2:water-on-earth:ocean}{Ocean}.
The largest is the Pacific \omterm{def:g2:water-on-earth:ocean}{Ocean}; the smallest is the Arctic \omterm{def:g2:water-on-earth:ocean}{Ocean}.
\end{solution}

\begin{solution}{exo:g3:continents-and-oceans:3}
The \omterm{def:g3:continents-and-oceans:poles}{North Pole} is at the top of a \omterm{def:g3:continents-and-oceans:globe}{globe}; the \omterm{def:g3:continents-and-oceans:poles}{South Pole} is at the bottom;
the \omterm{def:g3:continents-and-oceans:equator}{equator} goes all the way round the middle.
\end{solution}

\begin{solution}{exo:g3:continents-and-oceans:4}
Europe: northern. Australia: southern. Antarctica: southern. North America:
northern.
\end{solution}

\begin{solution}{exo:g3:continents-and-oceans:5}
The Arctic \omterm{def:g2:water-on-earth:ocean}{Ocean} lies around the \omterm{def:g3:continents-and-oceans:poles}{North Pole}. The Southern \omterm{def:g2:water-on-earth:ocean}{Ocean} lies around
Antarctica.
\end{solution}

\begin{solution}{exo:g3:continents-and-oceans:6}
From Africa to South America: the Atlantic \omterm{def:g2:water-on-earth:ocean}{Ocean}. From Australia east to
South America: the Pacific \omterm{def:g2:water-on-earth:ocean}{Ocean}. From Africa to Australia: the Indian
\omterm{def:g2:water-on-earth:ocean}{Ocean}.
\end{solution}

\begin{solution}{exo:g3:continents-and-oceans:7}
Sam lives near the \omterm{def:g3:continents-and-oceans:equator}{equator}, where it is hot all the year round. Bo lives
near a \omterm{def:g3:continents-and-oceans:poles}{pole}, where it is cold all the year round.
\end{solution}

\begin{solution}{exo:g3:continents-and-oceans:8}
$10 - 7 = 3$: her finger lands on the land 3 times. This is what we can
expect: about 71 squares out of 100 of the Earth are water, so about 7 times
out of 10 the finger stops on the \omterm{def:g2:water-on-earth:ocean}{sea}. (With only 10 spins, she could also
get 6 or 8: the count changes from one game to the next.)
\end{solution}

\begin{solution}{exo:g3:continents-and-oceans:9}
To draw the \omterm{def:g3:continents-and-oceans:world-map}{world map}, the \omterm{def:g3:continents-and-oceans:globe}{globe} was cut along a line through the middle of
the Pacific \omterm{def:g2:water-on-earth:ocean}{Ocean} and laid flat. One half of the Pacific ends up on the
left edge of the \omterm{def:g2:drawing-to-map:map}{map}, the other half on the right edge. There is only one
Pacific \omterm{def:g2:water-on-earth:ocean}{Ocean} on the real Earth: the two edges of the \omterm{def:g2:drawing-to-map:map}{map} are the same
line on the \omterm{def:g3:continents-and-oceans:globe}{globe}.
\end{solution}

\begin{solution}{exo:g3:continents-and-oceans:10}
Tom is wrong. From the largest to the smallest, the \omterm{def:g3:continents-and-oceans:continent}{continents} are Asia,
Africa, North America, South America, Antarctica, Europe and Australia:
Antarctica is only the fifth, much smaller than Asia. The \omterm{def:g3:continents-and-oceans:world-map}{world map} fools
him because a flat \omterm{def:g2:drawing-to-map:map}{map} stretches the land near the \omterm{def:g3:continents-and-oceans:poles}{poles} a lot.
Antarctica, which lies around the \omterm{def:g3:continents-and-oceans:poles}{South Pole}, is stretched into a long band
across the whole bottom of the \omterm{def:g2:drawing-to-map:map}{map}. On the \omterm{def:g3:continents-and-oceans:globe}{globe} it is a round mass of
land, and much smaller than Asia.
\end{solution}
